\subsection*{4.2 Upper bounds}
For our main upper bound, we modify our algorithm recreating the result of [11] (and analogously [12]) to have length $4n-4$. In particular, our program will read all but two variables four times, while the last two variables will be read twice.
\begin{proof}[Proof of Theorem3]
First, we make an easy change to Theorem9 which allows us to achieve $4n-3$. Observe that in ToggleInput the order in which we add the inputs is irrelevant, and so consider ToggleMonomials where we reverse the order of toggling on Line6. Then notice that the last query on Line2 and the first query on Line6 are both made to $x_n$, and so we can merge these two layers along with our entire for loop (which reads no variables) into a single layer querying $x_n$. Moving to Algorithm1, this means we only query $x_n$ twice, and furthermore the last query on Line1 and the first query on Line3 are both made to $x_1$, and so again by merging these two queries along with Line2 we only query $x_1$ three times.
Now we will change our program so that $x_1$ is only read twice. Consider two new functions obtained by fixing the value of $x_1$, namely $f^0=f(0,x_2\ldots x_n)$ and $f^1=f(1,x_2\ldots x_n)$. Recall that we used the following polynomial to compute $f$, where $y_i=\tau_i^{\mathrm{in}}+x_i$:
\[q_f=\sum_{S,S'\subseteq[n]}c_{S,S'}\left(\prod_{i\in S}\tau_i^{\mathrm{in}}\right)\left(\prod_{i\in S'}y_i\right).\]
If we choose $b\in\{0,1\}$ and fix $\tau_1^{\mathrm{in}}=0$ and $y_1=x_1=b$, we get the following, which can be used to compute (since it is equal to) $f^b$:
\[q_{f^b}=\sum_{S,S'\subseteq[2..n]}(c_{S,S'}+b\cdot c_{S,S'\cup\{1\}})\left(\prod_{i\in S}\tau_i^{\mathrm{in}}\right)\left(\prod_{i\in S'}y_i\right).\]
We will use Algorithm1 to compute $q_{f^b}$, where $b=x_1$, by removing all reference to $x_1$ from ToggleInput and ToggleMonomials, and querying $x_1$ whenever we execute Lines2 or4 to determine whether to compute $q_{f^0}$ or $q_{f^1}$ in place of $q_f$. More specifically, ToggleInput will now only loop over $i=2\ldots n$, while ToggleMonomials will now only loop over $S\subseteq[2..n]$. Finally in Algorithm1 we change Lines2 and4 to
\begin{multline*}
R_{\mathrm{out}}\leftarrow R_{\mathrm{out}}\pm\sum_{S,S'\subseteq[2..n]}(c_{S,S'}+x_1\cdot c_{S,S'\cup\{1\}})\\
\cdot\left(\prod_{i\in S}R_i^{\mathrm{in}}\right)R_{S'}
\end{multline*}
where Line2 uses+ and Line4 uses−. Note that to execute these lines correctly, we will query $x_1$ and perform the corresponding instruction; thus we no longer ignore these two lines in calculating our program length.
By our earlier definition of $q_{f^b}$, this exactly computes $q_{f^b}$ for $x_1=b$ as claimed. As above we will reverse the order of the queries in ToggleInput the second time it is called in ToggleMonomials, which allows us to read $x_n$ only once per execution for a total of two reads. $x_1$ will be queried in Lines2 and4, and all other variables will be queried four times.
Since we no longer use register $R_1$, or $R_S$ when $1\in S$, the number of registers is reduced from the original $n+2^n$ in Theorem9 to $n-1+2^{n-1}$. The output register $R_{\mathrm{out}}$ is now catalytic, so we add one register for a total $m=2^{n+2^{n-1}}$.
\end{proof}
